\section{Appendix: Scope and Qualifications}
\label{sec:scope}

{\small
\textbf{Edition, formulary and occurrence.} The study treats the Mass
\latin{Dominica decima nona post Pentecosten}, II classis, of the
\work{Missale Romanum}, editio typica, Typis Polyglottis Vaticanis 1962,
pp.~411--412, nos.~1679--1688, in the universal 1962 calendar and no particular
calendar. Its occurrence on 4~October 2026, the omission of the feast of St
Francis, the single oration of each kind, the colour, the Credo and the Trinity
Preface rest on the Missal's own calendar and general rubrics (RG 12, 16, 91,
111~b, 127~b; RGMR 433, 434~b, 475~a, 494~b), read on the page images of the
typical edition. The general rubrics tie the Gloria to the Te Deum of the
Office (RGMR 431~a), and the Breviary rubric behind it, which the typical
edition of the Missal does not print, was not read. The optional external solemnity of the Holy
Rosary on the first Sunday of October rests on RGMR 358~b and 360; the Rosary
formulary it brings into up to two Masses is not studied here. No diocesan,
titular, dedication or institute calendar is applied.

\textbf{Text control.} Every Latin form of the ten proper elements was read on
page images of the CMAA facsimile of the typical edition and collated with the
Pustet Missal of 1862, in which every element's wording already stands; the
collation, the conclusion forms and the comparison with the Clementine Vulgate
are recorded in \texttt{propers/verified.md}. The Introit's antiphon was not
traced to a scriptural or liturgical source, and the Latin source of the
Gradual's and Offertory's added \latin{Domine} and of the Offertory's future
tenses was not identified: the Vulgate does not account for them, and no Roman
or Old Latin psalter and no chant manuscript was consulted.

\textbf{English.} The scriptural English is the Douay--Rheims in Challoner's
revision at the canonical verses. Where the Missal departs from the Vulgate
(the Epistle's opening, the Gradual's \latin{Domine}, the titles left out of the
Introit's verse and the Alleluia, the Gospel's opening, and the Offertory's
\latin{Domine} and future tenses) that English does not render the Missal's
Latin, and no English has been supplied in its place. The English of the
Introit's antiphon and of the three orations is the anonymous translation
printed by Cummiskey in Philadelphia in 1861, read in the library's tracked
transcription of those pages, which was checked against the page images when it
was registered. Neither English is an approved liturgical translation of the
1962 Missal. The Trinity Preface is cited and not reproduced. The English of
St Augustine's sermons and of his expositions of Pss~77, 104, 118, 137 and 140,
of St John Chrysostom on St Matthew and on Ephesians, and of St Irenaeus is that
of the nineteenth-century Nicene and Post-Nicene and Ante-Nicene Fathers.
English renderings of the Latin of St Gregory, St Jerome, St Hilary, St Thomas
Aquinas, and St Augustine's \work{Quaestiones evangeliorum} and exposition of
Ps~104 that are not drawn from a named translation are glosses of the Latin
printed beside them.

\textbf{Reception.} Each appointed passage was searched in the Latin and Greek
Fathers and the later Doctors, saints and liturgical commentators available to
this production, and the witnesses read are those named in the References at
their exact loci. St Jerome's commentaries on St Matthew and on Ephesians, St
Hilary's on St Matthew and on Pss~118 and 140, and St Augustine's
\work{Quaestiones evangeliorum} were read on page images of Migne's
\work{Patrologia Latina}; St Gregory's homily in the Latin Wikisource
transcription of Migne's text, which was not collated with the printed
columns; St Thomas Aquinas on Ephesians on the page images of the Liège
edition of 1857, and on St Matthew only in the text layer of the Venice
edition of 1745, from which he is summarised and not quoted. The Greek of St
John Chrysostom on Pss~137 and 140 and of the \work{Expositiones in Psalmos}
printed under St Athanasius's name was read in the text layers of Migne's
\work{Patrologia Graeca} and is described, not quoted. The medieval liturgical
commentators Honorius Augustodunensis and Sicard of Cremona, and with them
Rupert of Deutz, William Durandus and Berno of Reichenau, expounded a Mass
that shared this formulary's Epistle and some of its chants but read another
Gospel, and
Honorius and Sicard are cited only for the elements that Mass shares with this
one. Bl.\ Ildefonso Schuster and the continuation of \work{The Liturgical Year}
comment on this formulary as the 1962 Missal gives it. Among the witnesses not
reached are St Ambrose on Ps~118, Cassiodorus and Theodoret on the psalms, the
book of Origen's commentary on St Matthew that treats this parable, the
\work{Catena aurea} on St Matthew, St Anthony of Padua's sermon on this Gospel,
and St Alphonsus Liguori's Sunday sermons.

\textbf{Chronology.} The Date cells in Scriptural Date and Location come from
the Scripture chronology corpus, through the record generated for this
formulary under its default profile. That record gives the five psalms only the
modern critical boundary for the Psalter, which it takes from the introduction
to the Psalms in the \work{New American Bible, Revised Edition} (summarised,
not quoted), and holds no traditional date or era for any of them, whatever
their titles name. For the Gospel it holds five dates from the \work{Catholic
Encyclopedia}'s article on St Matthew, which reports Eusebius's tradition that
Matthew wrote his Gospel in Hebrew when he left Palestine, and Durand's date in
``The New Testament'' for the Aramaic original, of whose rendering into Greek
Durand says that nothing definite is known. For the Epistle it holds Ladeuze's
range, with the place of the captivity, Caesarea or Rome, left open as a much
mooted question, and the year of Prat's table, which sets the letter in the
Roman captivity. The encyclopedia's articles do not agree on these dates, and
the record marks every one disputed. The Gospel's narrated event is undated. The record and its
sources are audited in \texttt{research/scope.md}.

\textbf{Rights.} The Latin of the formulary is published on the basis of its
public-domain antecedent, the Pustet Missal of 1862 (17 U.S.C. 103(b)). The
Douay--Rheims, the 1861 Cummiskey English, Migne's Latin and Greek texts, the
Nicene and Post-Nicene and Ante-Nicene Fathers, the Venice 1745 and Liège 1857
Aquinas, O'Sullivan's 1866 English of Bellarmine, the 1883 English of the
continuation of \work{The Liturgical Year} and the 1927 English of Schuster's
\work{Sacramentary} are in the public domain in the United States. St Gregory's
Latin is quoted from the Latin Wikisource transcription of the
\work{Homiliarum in Evangelia}, licensed under Creative Commons
Attribution-ShareAlike 3.0, with attribution to its contributors at
\url{https://la.wikisource.org/wiki/Homiliarum_in_Evangelia}. The New American
Bible introduction is summarised, not reproduced.

\textbf{Records.} The appointed-text collation is in \texttt{propers/verified.md};
the search, loci, disagreements and negative results in
\texttt{research/scope.md}; and the reasoning behind the three readings in
\texttt{research/interpretations.md}. The research records were reviewed
independently before this study was written.
\par}

\section*{References}
% A starred heading sets no mark, so the Appendix that ends above would go on
% heading the pages of the bibliography. The heading stays starred, because
% check-content-preflight finds the References section by that exact form; the
% mark is set explicitly through the leaf-local \runninghead, which the web
% converter's macro whitelist carries.
\runninghead{References}
\addcontentsline{toc}{section}{References}
\label{sec:references}

{\footnotesize
\subsection*{Liturgical books and Bibles}
\begin{itemize}
\item \work{Missale Romanum ex decreto Sacrosancti Concilii Tridentini restitutum, Summorum Pontificum cura recognitum}, editio typica (Typis Polyglottis Vaticanis, 1962), CMAA facsimile, \url{https://media.churchmusicassociation.org/pdf/missale62.pdf}: \work{Rubricae generales} 12, 16, 91, 111, 115, 127; \work{Rubricae generales Missalis} 358, 360, 427, 431, 433, 434, 475, 494; \work{Ordo Missae} nos.~1032, 1037, 1137; pp.~411--412 (nos.~1679--1688).
\item \work{Missale Romanum} (Ratisbon: Pustet, 1862), pp.~351--352, \latin{Dominica XIX. post Pentecosten}.
\item \work{The Roman Missal, Translated into the English Language for the Use of the Laity} (Philadelphia: Eugene Cummiskey, 1861), pp.~445--448 (XIX Sunday after Pentecost, Introit, Collect, Secret, Postcommunion).
\item The Holy Bible, Douay--Rheims version, revised by Bishop Richard Challoner (Project Gutenberg text): 1 Par 15:3; 16:1, 7--8; Ps 34:3; 77:1--8, 12--13, 49, 52--54, 67--72; 90:15; 104:1, 45; 118:1, 4--5, 19, 54; 137:1--4, 7; 138:16; 140:1--4; Mt 13:35; 21:1--18, 23, 43, 45; 22:1--15; 25:34--35; Rom 13:14; Eph 1:1; 3:1; 4:1, 22--32; 6:20; 1 Tim 1:5; 2:8.
\item \work{Biblia Sacra Vulgatae editionis} (Clementine Vulgate), eBible.org text: Ps 77:1; 104:1; 118:1, 4--5; 137:1, 7; 140:1--2; Mt 22:1--14; Eph 4:1, 23--28.
\end{itemize}

\subsection*{Fathers, Doctors and saints}
\begin{itemize}
\item St Augustine, \work{Sermones} 90 and 95 (Sermons XL and XLV), trans. in Nicene and Post-Nicene Fathers, 1st ser., vol.~6 (CCEL text): Sermo 90, 1, 4--6, 9--10; Sermo 95, 4--7.
\item St Augustine, \work{Enarrationes in Psalmos}, trans. in Nicene and Post-Nicene Fathers, 1st ser., vol.~8 (CCEL text): on Ps~77, 1--2; 104, 1; 118, Aleph 4--5; 137, 10--11; 140, 2--3. Latin on Ps~104, 1 and 34--36 in Migne's text (Latin Wikisource transcription).
\item St Augustine, \work{Quaestiones evangeliorum} I.31, in Migne, PL 35 (Paris, 1845), col.~1329.
\item St Gregory the Great, \work{Homiliae in Evangelia} 38, 1, 3--7, 9--14, Latin Wikisource transcription of Migne's text, CC BY-SA 3.0.
\item St Jerome, \work{Commentarii in Matthaeum} III, on Mt 22:1--14, and \work{Commentarii in Epistolam ad Ephesios} II, on Eph 4:23--28, in Migne, PL 26 (Paris, 1845), cols.~159--161 and 507--512.
\item St Hilary of Poitiers, \work{Commentarius in Matthaeum} 22, 3--7, and \work{Tractatus super Psalmos}, in Ps.~118, Aleph 11--12, and in Ps.~140, 3--4, in Migne, PL 9 (Paris, 1844), cols.~1042--1044, 508--509 and 825--826.
\item St John Chrysostom, \work{Homilies on Matthew} 69, Nicene and Post-Nicene Fathers, 1st ser., vol.~10 (CCEL text); \work{Homilies on Ephesians} 13--14, Nicene and Post-Nicene Fathers, 1st ser., vol.~13 (CCEL text); \work{Expositio in Psalmos}, in Ps.~137 and 140, Migne, PG 55 (text layer; described, not quoted).
\item St Irenaeus of Lyons, \work{Against Heresies} IV.36.5--6, Ante-Nicene Fathers, vol.~1 (Buffalo, 1887), pp.~516--517.
\item \work{Expositiones in Psalmos} printed under the name of St Athanasius, in Ps.~77, 104 and 118, Migne, PG 27 (text layer; described, not quoted).
\item St Thomas Aquinas, \work{In Epistolam ad Ephesios} c.~4, lect.~7--9, in \work{In omnes D.~Pauli Apostoli epistolas commentaria}, vol.~2 (Liège: H.~Dessain, 1857), pp.~327--333.
\item St Thomas Aquinas, \work{Super Evangelium S.~Matthaei lectura} c.~22, in \work{Opera}, editio altera Veneta, vol.~3 (Venice, 1745), pp.~279--282; summarised only.
\item St Robert Bellarmine, \work{A Commentary on the Book of Psalms}, trans. J. O'Sullivan (1866), read in the eCatholic2000 transcription of that translation: on Pss~104:1; 118:4; 137:7; 140:2.
\item St Rabanus Maurus, \work{Commentariorum in Matthaeum libri VIII}, on Mt 22:8--10, in Migne, PL 107 (text layer).
\end{itemize}

\subsection*{Liturgical commentators}
\begin{itemize}
\item Bl.\ Ildefonso Schuster, \work{The Sacramentary (Liber Sacramentorum)}, vol.~3 (London: Burns Oates \& Washbourne, 1927), pp.~171--174.
\item The continuation of Prosper Guéranger, \work{The Liturgical Year} (Dom Lucien Fromage), \work{The Time after Pentecost}, vol.~2 (series vol.~11), trans. Laurence Shepherd (Dublin: James Duffy and Sons, 1883), pp.~425--433.
\item Honorius Augustodunensis, \work{Gemma animae} IV.86, in Migne, PL 172 (text layer).
\item Sicard of Cremona, \work{Mitrale} VIII.19, in Migne, PL 213 (text layer).
\end{itemize}

\subsection*{Chronology sources}
\begin{itemize}
\item \work{The Catholic Encyclopedia} (New York, 1907--1912): ``Epistle to the Ephesians'' (vol.~5, 1909); ``Gospel of St.\ Matthew'' (vol.~10, 1911); ``St.\ Paul'' (vol.~11, 1911); ``The New Testament'' (vol.~14, 1912).
\item United States Conference of Catholic Bishops, \work{New American Bible, Revised Edition}, introduction to the Psalms, on their dating; summarised, not reproduced.
\end{itemize}
\par}
